\documentclass[parskip=full]{scrreprt}

\RedeclareSectionCommand[pagestyle=plain,afterskip=10pt plus 1pt]{chapter}
\setkomafont{chapter}{\mdseries\headingfont\fontsize{40}{40}\selectfont\color{black!80}}
\setkomafont{pageheadfoot}{\normalsize}

\def\pnumbox#1{#1\hspace*{8cm}}
\def\gobble#1{}
\DeclareTOCStyleEntry[
  indent=0pt,
  beforeskip=0pt,
  entryformat=\itshape,
  entrynumberformat=\textcolor{oldred},
  numwidth=2em,
  linefill=\hfill,
  pagenumberbox=\pnumbox,
  pagenumberformat=\itshape
]{tocline}{part}

% use commented values if there are no parts
\DeclareTOCStyleEntry[
  indent=0pt,
  beforeskip=0pt,
  entrynumberformat=\textcolor{oldred},
  numwidth=2em,
  linefill=\hfill,
  pagenumberbox=\pnumbox
]{tocline}{section}

\DeclareTOCStyleEntry[
  indent=0pt,
  beforeskip=-\parskip,
  entrynumberformat=\gobble,
  entryformat=\ltseries,
  numwidth=2em,
  linefill=\hfill,
  pagenumberbox=\pnumbox,
  pagenumberformat=\ltseries
]{tocline}{subsection}

\usepackage{polyglossia}
\setdefaultlanguage{english}

\usepackage{fontspec}

\setmainfont{Source Sans Pro}[
  ItalicFont = Source Sans Pro Italic,
  BoldFont = Source Sans Pro Bold,
  BoldItalicFont = Source Sans Pro Bold Italic,
  FontFace = {lt}{n}{Source Sans Pro Light},
  FontFace = {lt}{it}{Source Sans Pro Light Italic},
  FontFace = {sb}{n}{Source Sans Pro Semibold},
  FontFace = {sb}{it}{Source Sans Pro Semibold Italic},
  Numbers = Proportional,
  Ligatures = Common
]
\DeclareRobustCommand{\ltseries}{\fontseries{lt}\selectfont}
\DeclareRobustCommand{\sbseries}{\fontseries{sb}\selectfont}
\DeclareTextFontCommand{\textlt}{\ltseries}
\DeclareTextFontCommand{\textsb}{\sbseries}

\newfontfamily\headingfont{Fredericka the Great}

\usepackage[pass]{geometry}
\newgeometry{twoside,inner=20mm,outer=40mm,top=20mm,bottom=40mm}

\usepackage{url}
\urlstyle{same}

\usepackage{microtype}
\microtypesetup{verbose=silent}

\usepackage{booktabs,array,longtable}
\newcolumntype{L}[1]{>{\raggedright\let\newline\\\arraybackslash\hspace{0pt}}p{#1}}

\usepackage{graphicx}

\usepackage{xcolor}
\definecolor{oldred}{rgb}{.8313,0,0}

\usepackage{pdfpages}

\usepackage{scrlayer-scrpage}
\pagestyle{scrheadings}
\clearscrheadfoot
\cfoot[\thepage]{\thepage}
\pagenumbering{roman}

\usepackage{enumitem}
\setlist[description]{%
	labelindent=2em,%
	labelwidth=6em,%
	leftmargin=8em,%
	labelsep=0pt,
	first=\ltseries,%
	font=\normalfont\itshape\ltseries%
}
\def\lyrefitem#1{\item[\lyref{#1}]}


\makeatletter

\def\@firstofthree#1#2#3{#1}
\def\@secondofthree#1#2#3{#2}
\def\@thirdofthree#1#2#3{#3}
\def\@firstoffour#1#2#3#4{#1}
\def\@secondoffour#1#2#3#4{#2}
\def\@thirdoffour#1#2#3#4{#3}
\def\@fourthoffour#1#2#3#4{#4}
\def\Dotfill{\leavevmode\cleaders\hb@xt@ .75em{\hss .\hss }\hfill \kern \z@}

\def\lyrefnumber#1{\expandafter\@setref\csname r@#1\endcsname\@firstofthree{#1}}
\def\lyreftitle#1{\expandafter\@setref\csname r@#1\endcsname\@secondofthree{#1}}
\def\lyrefpage#1{\expandafter\@setref\csname r@#1\endcsname\@thirdofthree{#1}}

\def\lyrefgenrenumber#1{\expandafter\@setref\csname r@#1\endcsname\@firstoffour{#1}}
\def\lyrefgenregenre#1{\expandafter\@setref\csname r@#1\endcsname\@secondoffour{#1}}
\def\lyrefgenretitle#1{\expandafter\@setref\csname r@#1\endcsname\@thirdoffour{#1}}
\def\lyrefgenrepage#1{\expandafter\@setref\csname r@#1\endcsname\@fourthoffour{#1}}

\def\lyrefpart#1{%
	\begingroup%
	\makebox[0pt][l]{\sbseries\color{oldred}\lyrefnumber{#1}}\hspace*{2em}%
	\makebox[0pt][l]{\sbseries\lyreftitle{#1}}\hspace*{6em}%
	\hfill\sbseries\lyrefpage{#1}%
	\endgroup%
}
\def\lyrefsection#1{%
	\begingroup%
	\makebox[0pt][l]{\color{oldred}\lyrefgenrenumber{#1}}\hspace*{2em}%
	\makebox[0pt][l]{\ltseries\lyrefgenregenre{#1}}\hspace*{6em}%
	\lyrefgenretitle{#1}\Dotfill\lyrefgenrepage{#1}%
	\endgroup%
}
\InputIfFileExists{../out/lilypond.ref}{}{\InputIfFileExists{../lilypond.ref}{}{}}


\newcommand\fancytitlehead{
	\headingfont%
	\fontsize{80}{80}\selectfont\textcolor{black!80}{\@ifundefined{@shortname}{\@lastname}{\@shortname}.}\\[15pt]%
	\fontsize{50}{50}\selectfont\@ifundefined{@shorttitle}{\@title}{\@shorttitle}.%
}


\def\firstname#1{\def\@firstname{#1}}
\def\lastname#1{\def\@lastname{#1}}
\def\shortname#1{\def\@shortname{#1}}
\def\shorttitle#1{\def\@shorttitle{#1}}
\def\namesuffix#1{\def\@namesuffix{#1}}
\def\instrumentation#1{\def\@instrumentation{#1}}
\def\parts#1{\def\@parts{#1}}

\firstname{\relax}
\lastname{\relax}
\namesuffix{\relax}
\instrumentation{\relax}
\parts{\relax}


\def\maketitle{%
\begin{titlepage}%
	\Large%
	{\@titlehead}%
	\vfill%
	{\strut\@firstname}\\%
	{\sbseries\color{oldred}\strut\@lastname}\\%
	{\strut\@namesuffix}%
	\vfill%
	{\sbseries\@title}\\%
	{\@subtitle}\\[\baselineskip]%
	{\itshape\@instrumentation}%
	\vfill%
	{\itshape\@parts}\hspace*{\fill}\raisebox{0pt}[0pt][0pt]{\includegraphics{ees_logo}}%
\end{titlepage}%
}


\newif\iftemplate\templatetrue
\newif\ifprintreport\printreportfalse
\directlua{
scores = {
  ob1 = "Oboe I",
  ob2 = "Oboe II",
  ottoni = "Clarino I, II in C\string\\newline Timpani in C–G",
  vl1 = "Violino I",
  vl2 = "Violino II",
  vla = "Viola",
  coro = "Coro",
  org = "Organo",
  b = "Bassi",
  full_score = "Full Score"
}

last_arg = arg[\string#arg]
texio.write("Last argument: " .. last_arg)
if not (scores[last_arg] == nil) then
  tex.print("\string\\def\string\\lypdfname{" .. last_arg .. "}")
  tex.print("\string\\parts{" .. scores[last_arg] .. "}")
  if (last_arg == "full_score") then
    tex.print("\string\\printreporttrue")
  end
end
}


\@ifundefined{lypdfname}{%
  \templatefalse
  \printreporttrue
  \parts{Draft}
}{\templatetrue}

\makeatother






\begin{document}
\frenchspacing

\titlehead{\fancytitlehead}
\firstname{Johann Heinrich}
\lastname{Rolle}
\title{Der leidende Jesus}
\subtitle{Passions-Cantate\\WacR III:7\\(D-B Mus.ms 18709/1)}
\instrumentation{S, A, T, B (solo), S, A, T, B (coro), 2 ob, 2 clno, timp, 2 vl, vla, b, org}
\maketitle


\thispagestyle{empty}

\vspace*{\fill}

\raisebox{-4mm}{\includegraphics{byncsaeu}}\hspace*{1em}Wolfgang Esser-Skala, 2020

© 2020 by Wolfgang Esser-Skala. This edition is licensed under the Creative Commons Attribution-NonCommercial-ShareAlike 4.0 International License. To view a copy of this license, visit \url{http://creativecommons.org/licenses/by-nc-sa/4.0/}. 

Music engraving by LilyPond 2.18.0 (\url{http://www.lilypond.org}).\\
Front matter typeset with Source Sans Pro and Fredericka the Great.

\textit{First version, October 2020}

\vspace*{2cm}

\ifprintreport
\chapter*{Critical Report.}

This edition bases upon a copy in the Königliche Bibliothek zu Berlin, which has been digitized by the Staatsbibliothek zu Berlin – Preußischer Kulturbesitz. The digital version of the manuscript is available at \url{http://resolver.staatsbibliothek-berlin.de/SBB000281F600000000} (siglum Mus.ms. 18709/1).

In general, this edition closely follows the manuscript. Any changes that were introduced by the editor are indicated by italic type (lyrics, dynamics and directions), parentheses (expressive marks and bass figures) or dashes (slurs and ties). Accidentals are used according to modern conventions. Asterisks denote changes that are clarified in the detailed remarks below.\footnote{Abbreviations: A, alto; B, bass; b, basses; clno, clarion; Ms, manuscript; ob, oboe; org, organ; r, rest; S, soprano; T,~tenor; timp, timpani; vl, violin; vla, viola.}

\bigskip

\begin{longtable}{lll L{10cm}}
	\toprule
	\itshape Mov. & \itshape Bar & \itshape Staff & \itshape Note \\
	\midrule \endhead
	1.3 & 20  & A   & 1st quarter in Ms: g′4 \\
	1.6 & 109 & org & bar in Ms: G2–r4 \\
	\bottomrule
\end{longtable}

Bass figures appear in the following movements (bars in parentheses): 1.1, 1.3, 1.5 (1). The remaining bass figures were added by the editor.

This edition has been compiled and checked with utmost diligence. Nevertheless, errors and mistakes cannot be totally excluded. Please report any errors and mistakes to \url{wolfgang@esser-skala.at} or create an issue or pull request on the edition’s GitHub page \url{https://github.com/skafdasschaf/rolle-der-leidende-jesus-WacR-III-7}. Your help will be greatly appreciated.

\bigskip
\textit{Salzburg, October 2020\\
Wolfgang Esser-Skala}

\cleardoublepage
\chapter*{Contents.}

\lyrefpart{erstertheil}

\lyrefsection{omeineseel}

\begin{description}
	\item[Coro]
	O, meine Seel, ermuntre dich,\\
	daß ganz auf Jesu Leiden sich\\
	voll Andacht Herz und Auge lenkt,\\
	und Himmelswonne dann dich tränkt.
	
	Anbetend folge nun ihm nach,\\
	von Hohn zu Hohn, von Schmach zu Schmach,\\
	bis zu dem Kreuz, daran er stirbt\\
	und Sündern Gottes Huld erwirbt.
	
	Dann sage: Jesu, du mein Heil,\\
	ich bin dein Eigenthum und Theil,\\
	wie du mein bist, so bleib ich dein,\\
	und will mich ganz dir ewig weyhn.
\end{description}

\lyrefsection{derherrneiget}

\begin{description}
	\item[Coro]
	Der Herr neiget den Himmel und fähret herab.\\
	Dunkel ist unter seinen Füßen,\\
	und in finstern Wolken ist er verborgen.
\end{description}

\lyrefsection{jerusalem}

\begin{description}
	\item[Alto]
	Jerusalem! Jerusalem!\\
	Welch eine Nacht senkt über dich\\
	vom Himmel sich hernieder!\\
	Kommt Gott, der zum Gericht erwacht,\\
	und fodert der Propheten Blut\\
	von deinen Händen wieder?\\
	Steigt er zum Weltgericht herab von seinem Thron,\\
	der Oelberg erbebet, es rauschet der Kidron,\\
	von fern her rauscht der Donnerton.\\
	Wo ist der größte der Propheten, dein Jesus?\\
	Fleh ihn an, dich zu verteten.\\
	Wo find ich ihn, den meine Seele liebt?\\
	Hör ich ihn dort betrübt\\
	aus jener Mitternacht her klagen?\\
	Er ists.\\
	Mit welchen dunkeln Plagen\\
	hat ihn sein Gott geschlagen!
\end{description}

\lyrefsection{gerechter}

\begin{description}
	\item[Alto]
	Gerechter! Welche Qual erschrecket\\
	den eingebeugten Heldenmuth!\\
	Du bist statt Schweiß mit Blut bedecket,\\
	und Thränen fließen in dein Blut.\\
	Ach, weint der Heiligste selbst Zähren,\\
	faßt Schrecken ihn und banger Schmerz,\\
	wer wird mir Sünder Trost gewähren,\\
	wer gießet Ruh in dieses Herz?\\
	Wie soll ich vor dem Richter stehn,\\
	vor dem mein Jesus will vergehn?
\end{description}

\lyrefsection{labyrinth}

\begin{description}
	\item[Basso]
	O Labyrinth! der der Natur gebot,\\
	vor dem der Sturmwind schwieg,\\
	der dir, o Tod, allmächtig, seinen Raub entrißen,\\
	liegt hier, von Finsternissen\\
	umschattet und von Angst gedrängt,\\
	und leidet, was die Seele nicht gedenkt,\\
	und jammert Todestöne.\\
	Ach! wer erhellt mir diese dunkle Sonne?\\
	Mein Jesus, meiner Seelen Freund,\\
	was ist es, daß dein Auge weint?\\
	Heb aus dem Staube dich.\\
	Ein Blick von dir, nur einer, lehre mich,\\
	dies Todesjammern, dies Zagen\\
	verstehn und zu ertragen.\\
	Er höret mich: sanft ist sein Blick\\
	und ruft das Leben mir zurück.\\
	Sein Auge spricht: Die Arbeit dieser Nacht\\
	hat mir der Menschen Schuld gemacht.
\end{description}

\lyrefsection{singtihr}

\begin{description}
	\item[Coro]
	Singt, ihr Himmel, Gott ist Liebe,\\
	wunderbarlich ist sein Rath.\\
	Sing, o Erde, Gott ist Liebe,\\
	bey der Thaten größten That.\\
	Ihr Geschlechter der Erlösten,\\
	dem, der leidet, uns zu trösten,\\
	laßt uns ewig dankbar seyn.\\
	Jedes Herz, das ihn verkennet,\\
	nicht für ihn und Tugend brennet,\\
	faß einst ewig diese Pein.
\end{description}




\cleardoublepage
\fi

\iftemplate
\includepdf[pages=-]{../out/\lypdfname.pdf}
\fi

\end{document}